% !TEX root = main.tex
Motivated by the prevalence of the AI-human interplay in high-stake operations, we propose a \emph{learning to defer} model with limited and time-varying human capacity. When the cost distribution of jobs is known, we propose $\textsc{BACID}$ which balances the idiosyncrasy loss avoided by admitting a job and the delay loss the admission could incur to other jobs. We show that $\textsc{BACID}$ achieves a near-optimal $O(\sqrt{KT})$ regret, with $K$ being the number of job types and $T$ being the number of jobs. When the cost distribution of jobs is unknown, we show that an optimism-only extension of $\textsc{BACID}$ fails to learn because of the \emph{selective sampling} nature of the system. That is, humans only see jobs admitted by the AI, while labels from humans affect AI's classification accuracy. To address this issue, we introduce label-driven admissions. Finally, we extend our algorithm to a contextual setting and derive a regret guarantee of $\tilde{O}(T^{2/3})$ without dependence on the number of types. Our type aggregation technique for a many-class queueing system and analysis for queueing-delayed feedback enable us to provide (to the best of our knowledge) the first online learning result in contextual queueing systems. Moreover, we test our algorithm on real data and find that it can substantially reduce the number of misclassifications of existing content moderation practice.

Our work opens up several interesting questions:
\begin{itemize}
\item First, we assume that reviewers produce perfect labels with uniform speeds. In practice, reviewers have different expertise and can have errors or heterogeneous review speeds. How can we design the scheduling algorithms when reviewers have skill-dependent non-perfect review qualities and speeds? How can we learn these quantities for new reviewers?
\item Second, although we allow for time-varying arrival patterns and human reviewing capacity, we assume stationary cost distributions of jobs for the learning problem. In practice the cost distributions can change over time due to, e.g., the update of the offline ML models in content moderation. It is thus useful to extend our work to a non-stationary learning setting. However, as noted in Appendix~\ref{app:sim-fail-to-learn}, typical heuristics like discounted UCB may worsen the selective sampling issue in Section~\ref{sec:opti-fail}. Therefore, new algorithmic ideas are needed to address selective sampling in non-stationary settings.
\item Third, the main focus of this paper is to minimize the cost of misclassified jobs. Content moderation, however, has multiple objectives: for example, the platform may want to maintain high precision in the decisions. In this case, the problem becomes a constrained optimization with a constraint on the number of erroneously removed legitimate content. Although our model can incorporate this consideration by setting a value for legitimate content (Remark~\ref{remark:value}), how to find the most suitable value of legitimate content in a learning, congested, and non-stationary setting is however unclear.
\item Fourth, this paper assumes \emph{exogenous} arrivals of jobs. In content moderation, the jobs are content created by users who may change their behavior according to the platform's decision. For example, malicious users may create more policy-violating content when the human review system is congested. Tackling such \emph{endogenous} job arrival patterns is an open direction that can further improve the efficiency of content moderation practice.
\item Finally, while we, and a long line of queueing literature, focus on the fluid benchmark (and the $w-$fluid benchmark in Appendix~\ref{app:w-fluid}) to handle non-stationary arrivals and human capacities, they are not necessarily the tightest benchmarks. Identifying a suitable benchmark for a non-stationary queueing system remains an open question.
\end{itemize}